\section{Introduction}

Atomic positions are a compact experimental description of local material structure. From these coordinates, researchers can quantify strain, polarization, lattice rotation, defect topology, interface roughness, and local symmetry breaking, and can then correlate these structural descriptors with electronic, optical, magnetic, and transport properties \cite{Kalinin2022,Lee2020,Ziatdinov2017,krivanek2010atom}. This link is especially important in materials where functionality arises from deviations from ideal crystallinity, including point defects, grain boundaries, heterointerfaces, moir\'e reconstructions, surfaces, and aperiodic order \cite{Jiang2019,Rhodes2019,Nuckolls2024}. Atomic-resolution high-angle annular dark-field scanning transmission electron microscopy (HAADF-STEM) can image these structural motifs directly \cite{nellist2000principles,pennycook2012seeing}. However, the scientific value of such images depends on converting intensity contrast into accurate atomic coordinates. Here, we introduce Blob-Net, a deep-learning framework for robust, symmetry-agnostic atomic-column localization in atomic-resolution images.

Many algorithms have been developed to locate atomic columns in STEM images. Thresholding provides a simple route: pixels above a selected intensity cutoff are segmented, and atom positions are estimated from the centroids of the resulting regions \cite{Otsu1979Threshold}. This approach is fast, but it is sensitive to noise, nonuniform background, thickness contrast, and scan artifacts. Blob-detection methods, including Laplacian-of-Gaussian filters, identify bright compact features by enhancing local extrema across position and scale \cite{Lindeberg1998FeatureDetection}. These methods can estimate both feature centers and approximate radii, but their derivative-based kernels amplify noise and often require image-specific choices of scale, threshold, and peak-separation parameters. Gaussian fitting can provide high-precision coordinates by fitting a parametric intensity model to each atomic column \cite{Nord2017Atomap,Nomura2024Bayesian}. Under favorable signal-to-noise ratio, sampling, background, and model-validity conditions, this approach can reach picometer-scale precision. Yet Gaussian fitting still depends on initialization, peak assignment, and local model assumptions, and it can degrade near overlapping columns, interfaces, defects, background gradients, and distorted or low-dose images. These classical approaches rely heavily on parameter choices and can remain sensitive to noise, even when parameter tuning is automated through reward-guided optimization~\cite{Barakati2026Rewards}.

Machine-learning approaches have increasingly been used to reduce this parameter sensitivity. AtomAI \cite{Ziatdinov2022}, AtomSegNet \cite{Lin2021}, and related frameworks use convolutional neural networks (CNNs), often U-Net-like architectures, for microscopy segmentation tasks. Specialized models extend these approaches to overlapping atomic columns in bilayer moir\'e materials, where Gomb-Net enables layer-resolved atom identification through a multi-branch U-Net \cite{GombNet2025}, as well as defect identification in transition metal dichalcogenides \cite{Lee2020}, grain-boundary and phase mapping \cite{Vasudevan2018}, and low-signal-to-noise atom localization \cite{Lin2023AtomIDNet,Eliasson2024Localization}. Atom-segmentation networks have also been benchmarked across common datasets \cite{Wei2023Benchmark}. As the field moves toward automation and real-time decision-making at the microscope, rapid CNN inference enables workflows in which predicted atomic coordinates guide subsequent actions such as site-selective spectroscopy or electron-beam manipulation~\cite{Pratiush2024SmartSTEM}. Together, these studies show that CNNs can be fast, accurate, and robust when the training dataset is well matched to the experimental data.

The remaining limitation is not the U-Net architecture itself, but the structural prior embedded in the training data. Most atom-finding demonstrations are trained on images representing a single crystalline lattice geometry, defined by square, hexagonal, or otherwise periodic atomic arrangements. After a model is trained on periodic data, its reliance on periodic spatial correlations can lead to erroneous predictions at aperiodic regions such as nanoparticle edges, grain boundaries, and other structural defects. This behavior is consistent with shortcut learning, in which a model exploits correlations that perform well within the training distribution but fail to transfer to different conditions~\cite{Geirhos2020ShortcutLearning}. Here, the shortcut manifests as predictions that continue a learned lattice beyond image-supported material boundaries. In principle, this limitation can be mitigated by training on a sufficiently large and balanced dataset spanning all relevant atomic configurations. However, this is rarely feasible, and assumptions that limit the training dataset are projected into the model.

Here, we use an alternative training paradigm inspired by the original use of U-Net for biological image segmentation, where target features are spatially heterogeneous and nonperiodic \cite{Ronneberger2015}. We train a U-Net to localize atomic-column-like features in aperiodic HAADF-STEM images, thus optimizing the model on local atomic contrast rather than crystal-specific spatial priors. We compare Blob-Net with geometric models (Square-Net and Hex-Net, trained on periodic square and hexagonal arrangements, respectively) on simulated datasets, edge structures, and experimental images. Finally, we examine how receptive field, feature size, and feature spacing govern transfer across crystalline systems, establishing the conditions under which a single model can generalize across diverse atomic structures.

\section{Results}

\subsection{Comparison of models across geometries}

We first demonstrate the effect of training data geometry on atom localization. To isolate this effect, we simulate three identical datasets of HAADF-STEM images with different atom placements: square, hexagonal, and random. Three identical U-Net models are trained independently on these datasets. We refer to the resulting models as Square-Net, Hex-Net, and Blob-Net, where Blob-Net denotes the network trained on the random dataset. We refer to Square-Net and Hex-Net collectively as geometric models. The name emphasizes that the model is trained to recognize independent image features, or ``blobs,'' rather than to infer atom positions from larger spatial correlations.

Two distinct patterns emerge from testing all models across all datasets. Firstly, each model performs best on the dataset matching its training distribution (Fig.~\ref{fig:Blob-Net-1}). Square-Net achieves the highest accuracy on square lattices, Hex-Net on hexagonal lattices, and Blob-Net on the random dataset. This behavior is evident in both the F1 score, which measures whether the model detects the correct atom positions, and in the root-mean-square localization error (RMSE), which measures the distance of predicted positions to the ground truth position. In each training-dataset case, the native model achieves both the highest F1 score and the narrowest RMSE distribution.

\begin{figure}[!htbp]
    \centering
    \includegraphics[width=\textwidth]{fig-1.png}
    \caption[Blob-Net learns atom localization without a crystallographic training prior]{Blob-Net learns atom localization without a crystallographic training prior. \textbf{a,} Example image from each of the three simulated HAADF-STEM datasets. \textbf{b,} Corresponding ground truths. \textbf{c,} Model predictions from the three networks (columns) on the three example images (rows). \textbf{d,} Scatter plots of predicted atom coordinates minus ground-truth coordinates, with performance metrics (F1 score for atom prediction correctness and root-mean-square error (RMSE) for predicted coordinate offset from ground truth).}

    \label{fig:Blob-Net-1}
\end{figure}

The second pattern is that Blob-Net generalizes well to all datasets despite never observing a periodic lattice during training. On both square and hexagonal test datasets, its F1 score remains within 0.8\% of the corresponding native model. Although the localization RMSE increases by approximately 50\%, the predicted atom positions remain within one pixel of the ground truth positions. Because atom-finding algorithms are commonly followed by sub-pixel coordinate refinement, Blob-Net provides an excellent initialization even when the crystal geometry differs from its training distribution. In contrast, Square-Net and Hex-Net perform substantially worse when evaluated on geometries different from those seen during training, particularly on the random dataset. This asymmetry suggests that training on periodic data encourages the network to encode higher-order geometry information in addition to local atomic contrast. Blob-Net, by contrast, is trained without any repeating lattice motif and therefore learns primarily to identify atomic-column-like features based on their local contrast alone. Thus, Blob-Net is best understood as a learned local-contrast detector for atomic-column-like features. It identifies atoms first; it does not assume that those atoms must belong to a particular lattice. This combines the structure independence of classical contrast-based atom-finding methods with the rapid, hyperparameter-agnostic inference of machine-learning methods.

\subsection{Performance on edge structures}

Interesting material physics often occurs at defects, edges, and interfaces---the same structures where atom finding is the most difficult. Next, we test Blob-Net performance on these types of structures by comparing to Square-Net and Hex-Net performance on simulated HAADF-STEM images of defect boundaries in three crystals: molybdenum disulfide, strontium titanate, and graphene (Fig.~\ref{fig:Blob-Net-2}). Two important patterns emerge: Blob-Net reduces false-positive and false-negative detections across all datasets, and its localization RMSE is lowest for off-lattice atomic positions. The geometric models rely on patterns to predict positions; this reliance can lead to better localization in ideal cases but can be detrimental at edges and defects.

\begin{figure}[!htbp]
    \centering
    \includegraphics[width=\textwidth]{fig-2.png}
    \caption[Blob-Net reduces false-positive detections in structurally complex regions]{Blob-Net reduces false-positive detections in structurally complex regions. \textbf{a,} Simulated high-angle annular dark-field scanning transmission electron microscopy (HAADF-STEM) images of artificial edge structures of $\mathrm{MoS_2}$ (top), $\mathrm{SrTiO_3}$ (middle), and graphene (bottom). \textbf{b,} Corresponding ground truths. \textbf{c,} Model predictions from Square-Net, Hex-Net, and Blob-Net highlight true positives (TP), false positives (FP), and false negatives (FN). \textbf{d,} Scatter plots of predicted atom coordinates minus ground-truth coordinates, with performance metrics (F1 score for atom prediction correctness and root-mean-square error (RMSE) for predicted coordinate offset from ground truth).}

    \label{fig:Blob-Net-2}
\end{figure}


The dominant failure mode of the lattice-trained models is the extrapolation of periodic structure beyond the edge structure, leading to false-positive identifications. This behavior is consistent with shortcut learning based on periodic spatial correlations. A model trained only on periodic lattices may exploit spatial correlations that favor continuation of the lattice beyond the supported material boundary, even when the local intensity disagrees. This can be very useful for finding dim atoms, but can also lead to false-positive identifications at edge structures, potentially producing an incorrect structural interpretation, like in the case of Hex-Net predictions on the $\mathrm{MoS_2}$ edge (Fig. ~\ref{fig:Blob-Net-2}c). 

This effect can be mitigated by including edge structures in the training data, but this presents another problem. Consider the $\mathrm{MoS_2}$ case (Fig.~\ref{fig:Blob-Net-2}, top row), where the edge can be Mo-terminated or S-terminated. It is also possible for the boundary to proceed in a zig-zag direction, armchair direction, or any arbitrary combination of the two. For grain boundaries, the complexities increase with differences in high-angle and low-angle boundary physics and the incorporation of all known defect structures. The training data is constrained to these complexities for a fully understood system, but for a system where the common edge structures are not well understood, the training data must include a representative sample of all possible boundary configurations. 

Blob-Net takes the alternative approach and trains on aperiodically placed atoms. During training, the network is not rewarded for expecting long-range lattice motifs to continue. Instead, predictions are driven solely by local contrast, allowing the model to terminate immediately at material boundaries without spurious lattice continuation. Furthermore, the omission of any choice of edge structures during training avoids bias toward a predefined set of edge configurations. This approach reduces false positives and false negatives at edges and improves atom-localization robustness in the structurally heterogeneous test cases examined here.

The first two examples use small positional perturbations with a standard deviation of approximately 0.1 pixel in the lattice interior. In contrast, the graphene edge uses much larger perturbations of approximately 2.5 pixels, creating an extreme departure from the ideal lattice. This discontinuity approaching the edge disrupts the periodic lattice and reduces errors associated with shortcut learning, thus drastically reducing the false positive rates of all geometric model predictions on this dataset. 

The lattice disturbance in the graphene dataset also serves to test for model robustness. In the first two datasets, the two geometric models have high false-positive and false-negative prediction rates, but for true-positive predictions, they have very tight localization (within 0.30 pixels (px)). However, when the atoms are disturbed from the ideal lattice position (graphene dataset), Blob-Net RMSE (0.19 px) scores the best of all test cases. The substantial decline in geometric-model performance is consistent with the atom spacing lying outside the training distribution. Because Blob-Net was trained on aperiodic features, it has a high localization of features that have variable spacing, and thus there is a performance increase in this example. This phenomenon indicates that including evenly spaced features into the random dataset could improve localization on the images of crystal lattices, but this was not tested explicitly.


\subsection{Experimental image cases}

Next, we evaluate Blob-Net performance on experimental HAADF-STEM images of four materials with increasing geometric disorder: pristine monolayer $\mathrm{MoS_2}$, a $\Sigma3$ $\{111\} \langle100\rangle$ coherent twin grain boundary in face-centered cubic aluminum (FCC-Al) viewed along the $\langle110\rangle$, a high-angle grain boundary in monolayer $\mathrm{WS_2}$, and an $\mathrm{Al_{72}Ni_{11}Co_{17}}$ quasicrystal (Fig.~\ref{fig:Blob-Net-3}). Because these experimental images have no ground truth, we compare the results of Blob-Net to the results of Hex-Net to highlight the differences between models with and without a learned geometric prior.

The $\mathrm{MoS_2}$ monolayer is the ordered reference case, having almost no  geometric disorder. The $\Sigma3$ boundary is also relatively highly ordered. The two models perform nearly identically on these two images, except for a few false negatives from Hex-Net on the $\mathrm{MoS_2}$ image. These results confirm that the Blob-Net is comparable to geometric networks not only on simulated data, but also in experimental data.

\begin{figure}[!htbp]
    \centering
    \includegraphics[width=\textwidth]{fig-3.png}
    \caption[Comparison of Blob-Net and Hex-Net predictions on experimental structures]{Comparison of Blob-Net and Hex-Net predictions on experimental structures. Columns show \textbf{a,} pristine monolayer $\mathrm{MoS_2}$, \textbf{b,} $\Sigma3$ coherent twin grain boundary in face-centered cubic aluminum (FCC-Al), \textbf{c,} high-angle grain boundary in monolayer $\mathrm{WS_2}$, and \textbf{d,} $\mathrm{Al_{72}Ni_{11}Co_{17}}$ quasicrystal. The top row shows the experimental high-angle annular dark-field scanning transmission electron microscopy (HAADF-STEM) images and the bottom row shows the same images overlaid with predicted atomic-column positions. The overlays reveal agreement and differences between models rather than independently verified detection accuracy, because experimental ground-truth coordinates are unavailable. Scale bars, 1~nm.}

    \label{fig:Blob-Net-3}
\end{figure}

Increasing disorder in the $\mathrm{WS_2}$ grain boundary image begins to distinguish the performance of the two models. Grain boundaries present a particular difficulty for geometric models because they necessarily involve at least two crystalline regions of different orientations and the boundary region itself. At the displayed boundary, Blob-Net identifies four columns missed by Hex-Net, whereas Hex-Net identifies one column missed by Blob-Net. These differences could lead to materially different interpretations of the grain-boundary structure. 

Performance improvements in both networks could theoretically be achieved by constructing a training dataset specifically representing this $\mathrm{WS_2}$ monolayer and simulated grain-boundary defect structures. However, such targeted training would implicitly encode assumptions about the grain-boundary configurations used in training into the weights and biases of the network. Because, the space of possible grain-boundary structures is enormous, governed by five macroscopic degrees of freedom and additional microscopic degrees of freedom \cite{Sutton1995Interfaces}, it is not feasible to simulate all possible boundaries comprehensively. While only a subset of these structures might be encountered by a microscopist, the choice of which boundary structures to include in the training set would bias the outputs of the network towards those structures. In contrast, Blob-Net achieves good prediction on these types of structures without any encoded structural bias.

Increasing disorder further to the $\mathrm{Al_{72}Ni_{11}Co_{17}}$ quasicrystal image tests Blob-Net performance on an aperiodic ordered case. Quasicrystals are not random; they have long-range quasiperiodic order, but they lack translational periodicity. Because there is no finite repeating unit cell, it is not practical to simulate every local atomic environment that may appear in the image. This complexity makes the training of a geometric model for quasicrystals infeasible, and the performance of Hex-Net drops off drastically in this example as a result. In contrast, Blob-Net detects atoms at nearly all visually apparent atomic-column positions in the displayed region, although their accuracy cannot be established without ground truth. This example illustrates the potential value of geometry-agnostic training for aperiodically ordered structures.


\subsection{Feature scale, receptive field, and transfer}

We next quantify the length scales over which Blob-Net is effective, expressed relative to its fixed 68~$\times$~68~px bottleneck receptive field (RF). The RF is the region of the input image capable of influencing a feature at a given network layer. Each pixel in the neural network is only informed of the features within its own receptive field, so the RF sets a limit on the scale of features that can be learned. For the geometric models, it is important that the RF be large enough to capture periodic features in the training images, limiting the range of feature sizes and spacings that can be learned. Blob-Net does not have this limitation, so here we test the performance across feature sizes and spacings relative to the fixed RF by varying the physical pixel size of the simulated images to find Blob-Net's scale-transfer limitations.

Experimental pixel size in HAADF-STEM imaging is set by the microscope magnification and the scan-coil step size, and is independent of the size of the features in the image. Thus, when pixel size changes, the spacing between features changes and the size of the features in pixel units also changes. Here we test Blob-Net predictions across feature spacing while keeping feature size constant (Fig.~\ref{fig4}b), and across feature size while keeping feature spacing constant (Fig.~\ref{fig4}c). Because the geometric models rely on patterns of atoms with relations to their neighbors, those models are bound to keep feature spacing smaller than the RF. However, Blob-Net retains useful performance for the tested feature spacings larger than the bottleneck RF (the case where the network can fundamentally only see one atom at a time) because it does never learns a periodic prior. 

\begin{figure}[!htbp]
    \centering
    \includegraphics[width=\textwidth]{fig-4.png}
    \caption[Blob-Net transfer depends on feature scale relative to its receptive field]{Blob-Net transfer depends on feature scale relative to its receptive field. \textbf{a,} Schematic showing how the receptive field of a single pixel expands deeper in a convolutional neural network according to kernel size, stride, and downsampling. \textbf{b,} Simulated high-angle annular dark-field scanning transmission electron microscopy (HAADF-STEM) image of atoms with different spacings and the corresponding Blob-Net prediction. \textbf{c,} Simulated HAADF-STEM image of atoms with different atomic numbers and the corresponding Blob-Net prediction. \textbf{d,} Localization metrics across physical pixel sizes (top) and the corresponding atom-size and spacing ratios to the fixed 68~$\times$~68~pixel (px) bottleneck receptive field (RF) (bottom). The vertical dashed line represents the pixel size used for generating the dataset and training Blob-Net. Smaller physical pixels (left) correspond to a ``zoomed-in'' view and larger physical pixels (right) correspond to a ``zoomed-out'' view. The receptive field remains fixed in pixel units throughout the sweep.}

    \label{fig4}
\end{figure}

To quantify this, Blob-Net is tested on images from the original random dataset that are resampled across pixel sizes (Fig.~\ref{fig4}d). The original dataset has a pixel size of 0.1062~\AA/pixel, and at this pixel size, the model achieves F1~=~0.9730 on the testing dataset. As pixel size changes, Blob-Net remains above F1~=~0.95 between roughly 0.5--1.5$\times$ the trained pixel size. To the left of the plot the image is "zoomed in," and approaches a single atom occupying 40 percent of the RF. In this extreme, the model F1 score stays at 0.59 but RMSE increases to 4.2~px; the model is finding about half of the atoms, but the predicted positions are shifted 4 pixels away from the ground truth on average (although the atoms are very large, roughly 27 pixels wide). To the right of the plot, the image is "zoomed out," and approaches a single atom occupying a small fraction of the RF (3-4 pixels). In this extreme, the F1 score drops to 0.05, but the RMSE remains below a pixel (0.50). The model is finding very few of the atoms (because they are so small), but the predicted positions are within half a pixel from the ground truth positions (although the atoms are very small, only 3 to 4 pixels wide).

These results show Blob-Net is robust to pixel size change, but fails when features and spacing either dominate the RF or begin to disappear into the background of the RF. This analysis provides a direct quantitative characterization of scale-dependent failure for this atom-localization model and is important for understanding the limits of transferability of similar segmentation models. Blob-Net shows a comparatively broad useful pixel-size range, while prior studies likewise report deterioration outside the training scale distribution~\cite{Wei2023Benchmark,Sytwu2024Generalization}.

\section{Conclusions}

Blob-Net demonstrates that robust atom localization does not require perfectly representative training data. Rather, by training a standard U-Net on randomly placed atomic features, we obtain a model that remains competitive with geometric models on periodic lattices while achieving better performance at edges and defects in crystals, where interesting physics usually occur. Aperiodic training reduces the continuation of model predictions beyond image-supported material boundaries and supports localization where atomic positions depart substantially from an ideal lattice. Experimental comparisons further illustrate its utility across ordered crystals, grain boundaries, and quasiperiodic structures. The central contribution is therefore a training strategy rather than a new architecture: learning to identify atomic features without requiring their arrangement to match a predefined structural model.

This geometric flexibility does not imply unrestricted transfer. Performance remains sensitive to image contrast and feature scale, and the pixel-size sweep identifies a practical operating range beyond which detection and localization deteriorate. These results emphasize that removing a periodic prior and matching the imaging conditions are complementary requirements for reliable atom finding. Structure-specific training may improve performance on an individual image, but it also introduces assumptions about the configurations being measured. Blob-Net offers an alternative when those configurations are unknown or too diverse to simulate comprehensively, providing initial atomic coordinates for subsequent refinement and structural analysis. Its greatest value lies in enabling the study of unfamiliar atomic arrangements without first prescribing what those arrangements should be.

\section{Methods}

Synthetic datasets were generated as normalized HAADF-STEM-like images with known atom-center coordinates and Gaussian target heatmaps. The square, hexagonal, and random datasets each used 512~$\times$~512~px images and train/validation/test splits of 1024/256/256 images. These datasets were matched in intensity, atom width, atom spacing, background, noise, blur, and edge-padding ranges so that the primary difference was geometry. Square and hexagonal datasets used periodic lattice generators with randomized rotations, lattice spacings, position jitter, and vacancy fractions. The random dataset sampled nonperiodic atom centers subject to a per-image minimum separation. All datasets included randomized atom intensity, atom width, constant background, linear gradients, smooth inhomogeneous background, low-frequency noise, Poisson shot noise, read noise, Gaussian blur, and off-frame edge padding.

Image contrast, noise, and sampling were matched to the training data, while the material boundaries were synthetically modified to introduce local disorder and remove long-range periodicity. These images therefore retain realistic imaging conditions while providing a controlled test of model performance at structurally complex edges.

The model was a single-output U-Net for heatmap regression \cite{Ronneberger2015}. The encoder and decoder used convolutional blocks with two $3 \times 3$ convolutions, batch normalization, and rectified linear unit (ReLU) activations. The filter widths were [32, 64, 128, 256], with dropout 0.2 at the bottleneck. The models trained on the square, hexagonal, and random datasets are referred to as Square-Net, Hex-Net, and Blob-Net. The training loss combined global mean-squared error and a peak-weighted mean-squared error. The peak term used target values greater than 0.1 and a peak boost of 5.0. Models were trained for up to 20 epochs with batch size 32, learning rate 0.001, and early stopping patience of 6 epochs.

For localization, model heatmaps were thresholded relative to the heatmap maximum, local maxima were detected with a minimum peak separation, and each peak was refined by a weighted centroid in a local window. Model predictions on synthetic images were matched to ground-truth coordinates by Hungarian assignment in pixel units. The cross-geometry manuscript comparison used a relative threshold of 0.35 and a maximum match distance of 3.0~px. Performance was summarized by precision, recall, F1 score, mean localization error, median localization error, and RMSE.

Experimental images were processed by applying the trained Blob-Net model to HAADF-STEM images and inspecting the resulting heatmaps and detected coordinates. Because experimental ground-truth coordinates are unavailable, these cases should be interpreted as qualitative transfer tests unless paired with manual annotations, crystallographic reference coordinates, or independent coordinate-refinement workflows.
